% AI-written, AI cross-reviewed; no human sign-off; not a claim that RH is proved
% Atlas unit D575. Local manuscript v3, item G.614, heading on printed PDF p. 310.
% GlobalTheory.GreenBridge.green_of_regular
% paper2/decomposition_tables.tex:5096–5101. Verbatim excerpt; SHA-256 210d70168a590d79e5353da1a72d4b28439d493cd3dabffa7f1e89e995de7b05.
% Original TeX uses the macros and cross-reference keys of the source manuscript.


% Start of atlas edition annotation. The following block is outside the original excerpt.
\paragraph{Status of named inputs in this reading.}
\noindent\nolinkurl{GlobalTheory.Published.GreenRegularCharacters} --- hypothesis at this head (\path{Development/AbelianVarietyTheory/PublishedTheorems/GaloisCohomology₂/GreenRegularCharacters.lean:13--20}); variant.\par
\noindent\nolinkurl{UpperTheory.LiteratureTheorems.green_regular_characters_of_literature} --- in-repo constant (\path{Development/AutomorphicArithmetic/ThetaCorrespondence/ThetaPairLifts/ThirdPlaceSeries/EveryCharacterOfLiteratureOfVWin.lean:399--400}); variant.\par
\noindent\nolinkurl{GlobalTheory.GreenBridge.green_of_regular} --- in-repo constant (\path{Development/AbelianVarietyTheory/Proofs/InertiaEigenvalues/SumRangeEqRoots.lean:225--304}); not compared.\par
% End of atlas edition annotation.

\begin{formalstatement}[D575]\label{decomp:D575}
Assume the regular-character classification for finite general linear groups. For every prime field size g, permitted character pairing L and degree d>0, index the cuspidal characters by monic irreducible polynomials of degree d with nonzero constant coefficient. The characters have the prescribed degrees and Frobenius-orbit value formula; the labels identify equal characters and exhaust the cuspidal ones.
\end{formalstatement}
{\small\raggedright\noindent\textsf{\bfseries Related reference}\quad \nolinkurl{GlobalTheory.GreenBridge.green_of_regular}\par\noindent \nolinkurl{Development/AbelianVarietyTheory/Proofs/InertiaEigenvalues/SumRangeEqRoots.lean:225--304}\par\noindent SHA-256 \nolinkurl{bf88f05c856ea17023e795b5f4927ca92a35971bd50f00ca4ed56ad8505d16f7}\par\noindent Dominated 280\quad Disjoint 0\quad terminal.\par}
Green gives the character parametrization. Changing from regular characters modulo Frobenius to the fixed polynomial roots and the selected pairing requires root, orbit and minimal-polynomial comparisons.
\hyperref[lit:S065]{Published sources and search S065}.
